
\documentclass[11pt,a4paper]{article}

\usepackage[utf8]{inputenc}
\usepackage[T1]{fontenc}
\usepackage[french]{babel}
\usepackage{lmodern}
\usepackage{microtype}
\usepackage[a4paper,margin=1.65cm,headheight=15pt]{geometry}
\usepackage{amsmath,amssymb,mathtools}
\usepackage{array,tabularx,multicol}
\usepackage{enumitem}
\usepackage[most]{tcolorbox}
\usepackage{xcolor}
\usepackage{tikz}
\usetikzlibrary{arrows.meta,calc}
\usepackage{pdflscape}
\usepackage{fancyhdr}
\usepackage{fancyvrb}
\usepackage{listings}
\usepackage{hyperref}

\definecolor{fondpage}{RGB}{247,248,250}
\definecolor{encre}{RGB}{31,35,40}
\definecolor{bleu}{RGB}{40,91,135}
\definecolor{vert}{RGB}{55,125,92}
\definecolor{orange}{RGB}{178,105,42}
\definecolor{violet}{RGB}{101,77,145}
\definecolor{rouge}{RGB}{170,72,72}
\definecolor{gris}{RGB}{86,91,99}
\definecolor{bleufonce}{RGB}{28,55,120}
\definecolor{bleuclair}{RGB}{238,243,255}
\definecolor{grisclair}{RGB}{245,245,245}
\definecolor{tableline}{RGB}{45,45,45}
\definecolor{softgray}{RGB}{246,247,249}
\definecolor{deepblue}{RGB}{25,62,115}
\definecolor{softblue}{RGB}{235,241,250}
\definecolor{accent}{RGB}{20,82,130}

\hypersetup{colorlinks=true,linkcolor=bleufonce,urlcolor=bleufonce,pdftitle={Parcours de révision -- Géométrie dans l'espace -- Terminale spécialité},pdfauthor={}}

\setlength{\parindent}{0pt}
\setlength{\parskip}{0.25em}
\renewcommand{\arraystretch}{1.12}
\setlength{\tabcolsep}{3pt}
\setlist[itemize]{leftmargin=1.25em,itemsep=0.15em,topsep=0.2em}
\setlist[enumerate,1]{label=\textbf{\arabic*.}, leftmargin=1.05cm, itemsep=0.3em, topsep=0.2em}
\setlist[enumerate,2]{label=\textbf{\alph*.}, leftmargin=0.85cm, itemsep=0.2em, topsep=0.15em}

\pagestyle{fancy}
\fancyhf{}
\lhead{}
\rhead{}
\cfoot{\thepage}

\newcommand{\R}{\mathbb{R}}
\newcommand{\N}{\mathbb{N}}
\newcommand{\vect}[1]{\overrightarrow{#1}}
\newcommand{\vu}{\vec u}
\newcommand{\vv}{\vec v}
\newcommand{\vw}{\vec w}
\newcommand{\vn}{\vec n}
\newcommand{\vd}{\vec d}
\newcommand{\norme}[1]{\left\lVert #1\right\rVert}
\newcommand{\ps}{\cdot}
\newcommand{\e}{\mathrm e}
\newcommand{\origine}[1]{ {\scriptsize\textcolor{bleufonce}{(site #1)}}}

\newcounter{question}
\newcounter{reponse}
\newenvironment{blocquestions}[1]{%
  \begin{tcolorbox}[
    enhanced,
    breakable,
    colback=bleuclair,
    colframe=bleufonce,
    coltitle=white,
    colbacktitle=bleufonce,
    title={\bfseries #1},
    fonttitle=\bfseries,
    boxrule=0.6pt,
    arc=2mm,
    left=2mm,
    right=2mm,
    top=2.5mm,
    bottom=1.5mm,
    toptitle=1.2mm,
    bottomtitle=1.2mm,
    lefttitle=2mm,
    righttitle=2mm,
    before upper=\vspace{0.7mm},
    before skip=3.5mm,
    after skip=3mm, fontupper=\large
  ]
  \large
  \begin{enumerate}[leftmargin=*,label=\textbf{\arabic*.},itemsep=0.85em,topsep=0.5em]
  \setcounter{enumi}{\value{question}}
}{%
  \setcounter{question}{\value{enumi}}
  \end{enumerate}
  \end{tcolorbox}
}
\newenvironment{blocreponses}[1]{%
  \begin{tcolorbox}[
    enhanced,
    breakable,
    colback=bleuclair,
    colframe=bleufonce,
    coltitle=white,
    colbacktitle=bleufonce,
    title={\bfseries #1},
    fonttitle=\bfseries,
    boxrule=0.6pt,
    arc=2mm,
    left=2mm,
    right=2mm,
    top=2.5mm,
    bottom=1.5mm,
    toptitle=1.2mm,
    bottomtitle=1.2mm,
    lefttitle=2mm,
    righttitle=2mm,
    before upper=\vspace{0.7mm},
    before skip=3.5mm,
    after skip=3mm, fontupper=\large
  ]
  \large
  \begin{enumerate}[leftmargin=*,label=\textbf{\arabic*.},itemsep=0.45em,topsep=0.5em]
  \setcounter{enumi}{\value{reponse}}
}{%
  \setcounter{reponse}{\value{enumi}}
  \end{enumerate}
  \end{tcolorbox}
}

\newtcolorbox{correctionbox}[1]{enhanced,breakable,colback=grisclair,colframe=black!55,boxrule=0.55pt,arc=2mm,left=2mm,right=2mm,top=1.5mm,bottom=1.5mm,title=\textbf{#1},coltitle=black,colbacktitle=black!12,before skip=2mm,after skip=2mm}
\newcommand{\titreSujet}[2]{%
  \subsection*{#1}
  \addcontentsline{toc}{subsection}{#1}
  \textit{Référence : #2}\par\medskip\hrule\medskip
}

% Styles de la carte mentale de cours
\setlist[itemize]{
    leftmargin=1.1em,
    itemsep=0.10em,
    topsep=0.15em,
    parsep=0pt,
    partopsep=0pt,
    label=\textcolor{black!65}{\raisebox{0.15ex}{\tiny$\blacktriangleright$}}
}
\tcbset{
    enhanced,
    arc=2.2mm,
    outer arc=2.2mm,
    boxrule=0.45pt,
    left=1.55mm,
    right=1.55mm,
    top=1.65mm,
    bottom=1.65mm,
    boxsep=0.8mm,
    before skip=1.6mm,
    after skip=1.6mm,
    fonttitle=\bfseries\small,
    coltitle=encre,
    before upper=\raggedright\fontsize{9.05}{10.9}\selectfont,
}
\newtcolorbox{fichebox}[3][]{
    colback=white,
    colframe=#2!72!black,
    colbacktitle=#2!15,
    coltitle=#2!55!black,
    borderline west={1mm}{0pt}{#2!78!black},
    title={#3},
    drop fuzzy shadow=black!4,
    #1
}
\newcommand{\tagcol}[2]{%
\begin{tcolorbox}[
    enhanced,
    colback=#1!11,
    colframe=#1!45!black,
    boxrule=0.42pt,
    arc=2.2mm,
    left=1mm,right=1mm,top=0.55mm,bottom=0.55mm,
    before skip=0mm,
    after skip=1mm,
    fontupper=\bfseries\small,
    colupper=#1!55!black
]
\centering #2
\end{tcolorbox}}
\newcommand{\titrepage}{%
\begin{tcolorbox}[
    colback=white,
    colframe=bleu!70!black,
    boxrule=0.65pt,
    arc=3mm,
    left=2.5mm,right=2.5mm,top=1.2mm,bottom=1.2mm,
    before skip=0mm,
    after skip=1.6mm,
    drop fuzzy shadow=black!8
]
{\Large\bfseries Parcours de révision -- Géométrie dans l'espace}
\end{tcolorbox}}

\tikzset{tabouter/.style={draw=tableline, line width=0.85pt},tabline/.style={draw=tableline, line width=0.55pt},forbidden/.style={draw=tableline, line width=0.95pt},vararrow/.style={-{Latex[length=2.7mm,width=2.0mm]}, line width=0.85pt, draw=accent},tabtext/.style={font=\small},tabmath/.style={font=\small},labelcell/.style={font=\small\bfseries, text=deepblue},pointvalue/.style={font=\small\bfseries},endvalue/.style={font=\small}}
\lstset{language=Python,basicstyle=\ttfamily\small,numbers=left,numberstyle=\tiny,frame=single,showstringspaces=false,columns=fullflexible,keepspaces=true}


\begin{document}
\begin{center}
{\LARGE\bfseries Parcours de révision -- Géométrie dans l'espace}\\[1mm]

\end{center}
\vspace{2mm}
\tableofcontents
\clearpage

\phantomsection
\addcontentsline{toc}{section}{Partie 1 -- Récapitulatif de cours}
\newgeometry{margin=8mm}
\pagestyle{empty}
\begin{landscape}
\pagecolor{fondpage}\color{encre}
\thispagestyle{empty}
\fontsize{8.55}{10.05}\selectfont
\titrepage

\noindent
\begin{minipage}[t]{0.318\linewidth}
\tagcol{bleu}{1. REPÉRAGE ET VECTEURS}

\begin{fichebox}{bleu}{Coordonnées utiles}
Si $A(x_A;y_A;z_A)$ et $B(x_B;y_B;z_B)$ :
\[
\vect{AB}(x_B-x_A\,;\,y_B-y_A\,;\,z_B-z_A).
\]
Si $\vu(x;y;z)$ :
\[
\norme{\vu}=\sqrt{x^2+y^2+z^2}.
\]
\end{fichebox}

\begin{fichebox}{vert}{Distance et milieu}
\[
AB=\sqrt{(x_B-x_A)^2+(y_B-y_A)^2+(z_B-z_A)^2}.
\]
\[
I\left(\frac{x_A+x_B}{2};\frac{y_A+y_B}{2};\frac{z_A+z_B}{2}\right).
\]
\end{fichebox}

\begin{fichebox}{orange}{Colinéarité}
\[
\vu \text{ et } \vv \text{ colinéaires}\iff \exists k\in\R,\ \vu=k\vv.
\]
Coordonnées non proportionnelles $\Rightarrow$ vecteurs non colinéaires.
\end{fichebox}

\begin{fichebox}{orange}{Coplanarité}
\[
\vu,\vv,\vw \text{ coplanaires}\iff \vw=a\vu+b\vv.
\]
Trois vecteurs non coplanaires sont linéairement indépendants et forment une base de l'espace.
\end{fichebox}

\end{minipage}
\hfill
\begin{minipage}[t]{0.318\linewidth}
\tagcol{vert}{2. DROITES ET PLANS}

\begin{fichebox}{vert}{Droite paramétrée}
Une droite passant par $A(x_A;y_A;z_A)$ et de vecteur directeur $\vu(a;b;c)$ admet une représentation :
\[
\begin{cases}
 x=x_A+ta,\\[-0.5mm]
 y=y_A+tb,\\[-0.5mm]
 z=z_A+tc,
\end{cases}
\quad t\in\R.
\]
\end{fichebox}

\begin{fichebox}{vert}{Plan et vecteur normal}
Si $P$ admet pour vecteur normal $\vn(a;b;c)$, alors :
\[
P: ax+by+cz+d=0.
\]
Réciproquement, dans $ax+by+cz+d=0$, un vecteur normal est :
\[
\vn(a;b;c).
\]
\end{fichebox}

\begin{fichebox}{violet}{Équation de plan : réflexe}
Une équation est de la forme :
\[
ax+by+cz+d=0.
\]
Comme un point $A$ appartient au plan, ses coordonnées vérifient l'équation : cela permet de trouver $d$.
\end{fichebox}

\begin{fichebox}{bleu}{Positions relatives}
\begin{tabularx}{\linewidth}{>{\bfseries}l X}
Deux droites & sécantes, parallèles, confondues ou non coplanaires.\\
Deux plans & sécants selon une droite, parallèles ou confondus.\\
Droite/plan & sécants, parallèles ou droite incluse dans le plan.
\end{tabularx}
\end{fichebox}

\end{minipage}
\hfill
\begin{minipage}[t]{0.318\linewidth}
\tagcol{violet}{3. PRODUIT SCALAIRE, ORTHOGONALITÉ, DISTANCES}

\begin{fichebox}{vert}{Produit scalaire}
Dans un repère orthonormé, si $\vu(x;y;z)$ et $\vv(x';y';z')$ :
\[
\vu\ps\vv=xx'+yy'+zz'.
\]
Avec les normes :
\[
\vu\ps\vv=\norme{\vu}\,\norme{\vv}\cos(\theta).
\]
\end{fichebox}

\begin{fichebox}{vert}{Orthogonalité}
\[
\vu\perp \vv \iff \vu\ps\vv=0.
\]
Deux droites sont orthogonales si leurs vecteurs directeurs sont orthogonaux.

\medskip
Elles sont perpendiculaires si elles sont orthogonales et sécantes.
\end{fichebox}

\begin{fichebox}{violet}{Droite orthogonale à un plan}
Pour montrer qu'une droite est orthogonale à un plan, on montre qu'un vecteur directeur de la droite est orthogonal à \textbf{deux vecteurs non colinéaires} du plan.
\[
\vd\ps\vu=0 \quad\text{et}\quad \vd\ps\vv=0.
\]
Alors $\vd$ est un vecteur normal au plan.
\end{fichebox}

\begin{fichebox}{orange}{Projeté et distance}
Le projeté orthogonal est le point le plus proche.

Si $H$ est le projeté de $A$ sur une droite $d$ :
\[
H\in d \quad\text{et}\quad (AH)\perp d,
\qquad d(A,d)=AH.
\]
Si $H$ est le projeté de $A$ sur un plan $P$ :
\[
H\in P \quad\text{et}\quad (AH)\perp P,
\qquad d(A,P)=AH.
\]
\end{fichebox}

\end{minipage}
\end{landscape}

\restoregeometry
\pagestyle{fancy}
\nopagecolor
\color{black}
\setlist[itemize]{leftmargin=1.25em,itemsep=0.15em,topsep=0.2em,label=\textbullet}
\setlist[enumerate,1]{label=\textbf{\arabic*.}, leftmargin=1.05cm, itemsep=0.3em, topsep=0.2em}
\setlist[enumerate,2]{label=\textbf{\alph*.}, leftmargin=0.85cm, itemsep=0.2em, topsep=0.15em}
\clearpage

\section*{Partie 2 -- Questions flashs}
\addcontentsline{toc}{section}{Partie 2 -- Questions flashs}
\setcounter{question}{0}
\begin{blocquestions}{A. Repères, vecteurs, droites et plans}
  \item Dans l'espace, comment calcule-t-on les coordonnées de $\overrightarrow{AB}$ ? \origine{170}
  \item Comment calcule-t-on les coordonnées du milieu de $[AB]$ ? \origine{171}
  \item Comment montre-t-on que trois points $A$, $B$, $C$ sont alignés ? \origine{172}
  \item Que signifie « deux vecteurs sont colinéaires » ? \origine{173}
  \item Qu'est-ce qu'un vecteur directeur d'une droite ? \origine{174}
  \item Qu'est-ce qu'une représentation paramétrique d'une droite ? \origine{175}
  \item Dans une représentation paramétrique, quel rôle joue le paramètre ? \origine{176}
  \item Comment reconnaît-on qu'un point appartient à une droite paramétrée ? \origine{177}
  \item Qu'est-ce qu'un vecteur normal à un plan ? \origine{178}
  \item Si $\vec n=(a,b,c)$ est normal à un plan, quelle forme a une équation cartésienne de ce plan ? \origine{179}
  \item Comment vérifie-t-on qu'un point appartient à un plan d'équation $ax+by+cz+d=0$ ? \origine{180}
  \item Que fait-on souvent pour déterminer l'intersection de deux plans dont on connaît les équations ? \origine{181}
  \item Quand deux plans sont-ils parallèles ? \origine{182}
  \item Quand une droite est-elle parallèle à un plan ? \origine{183}
  \item Que signifie que quatre points sont coplanaires ? \origine{184}
  \item Comment peut-on montrer que quatre points $A$, $B$, $C$, $D$ ne sont pas coplanaires ? \origine{185}
\end{blocquestions}

\begin{blocquestions}{B. Distances, produit scalaire et normes}
  \item Quelle est la formule du volume d'une pyramide ? \origine{186}
  \item Dans un repère orthonormé, quelle est la formule du produit scalaire de $\vec u(x,y,z)$ et $\vec v(x',y',z')$ ? \origine{187}
  \item Que vaut $\vec u\cdot \vec u$ ? \origine{188}
  \item Quand deux vecteurs non nuls sont-ils orthogonaux ? \origine{189}
  \item Quelle formule relie produit scalaire et angle ? \origine{190}
  \item Comment calcule-t-on la norme d'un vecteur de coordonnées $(x,y,z)$ ? \origine{191}
  \item Comment calcule-t-on la distance entre deux points $A$ et $B$ dans un repère orthonormé ? \origine{192}
  \item Quand une droite est-elle orthogonale à un plan ? \origine{193}
  \item Qu'est-ce que le projeté orthogonal d'un point sur un plan ? \origine{194}
  \item Si $H$ est le projeté orthogonal d'un point $A$ sur un plan $(P)$, que représente la longueur $AH$ ? \origine{195}
  \item Quel vecteur directeur lit-on sur $x=1+2t$, $y=3-t$, $z=4+5t$ ? \origine{250}
  \item Quel est le vecteur normal du plan $2x-y+3z-4=0$ ? \origine{251}
  \item Si deux vecteurs non nuls vérifient $\vec u\cdot\vec v=0$, que peut-on conclure ? \origine{252}
\end{blocquestions}

\begin{blocquestions}{C. Définitions et rédaction}
  \item Qu'appelle-t-on deux droites coplanaires ? \origine{295}
  \item Qu'appelle-t-on une équation cartésienne d'un plan ? \origine{316}
  \item Qu'appelle-t-on la norme d'un vecteur ? \origine{317}
  \item Dans un exercice de géométrie de l'espace, un élève écrit : « L'équation de $\vec u$ est $(1;2;3)$. » Comment formuler proprement la réponse pour avoir tous les points ? \origine{320}
\end{blocquestions}

\begin{blocquestions}{D. Calculs directs}
  \item Quelles sont les coordonnées de $\overrightarrow{AB}$ si $A(1;2;3)$ et $B(4;0;5)$ ? \origine{384}
  \item Que vaut le produit scalaire de $(1;0;2)$ et $(3;-1;1)$ ? \origine{386}
  \item Quel critère de produit scalaire permet de montrer qu'un triangle $ABC$ est rectangle en $A$ ? \origine{537}
  \item Quel est l'ensemble des points $M$ tels que $\overrightarrow{MA}\cdot\overrightarrow{MB}=0$ ? \origine{538}
  \item Comment détermine-t-on une équation cartésienne d'une droite connaissant un point et un vecteur normal ? \origine{539}
\end{blocquestions}
\clearpage
\section*{Partie 3 -- Sujets de bac}
\addcontentsline{toc}{section}{Partie 3 -- Sujets de bac}
\titreSujet{Sujet 1 -- Vrai / faux dans l'espace}{Amérique du Nord, 21 mai 2025, jour 1, exercice 3}

Pour chacune des affirmations suivantes, indiquer si elle est vraie ou fausse. Chaque réponse doit être justifiée. Une réponse non justifiée ne rapporte aucun point.

L'espace est rapporté à un repère orthonormé $(O;\vec{i},\vec{j},\vec{k})$.
On considère la droite $(d)$ dont une représentation paramétrique est :
\[
\begin{cases}
 x=3-2t,\\
 y=-1,\\
 z=2-6t,
\end{cases}\qquad t\in\mathbb{R}.
\]
On considère également les points suivants :
\[
A(3;-3;-2),\qquad B(5;-4;-1),
\]
$C$ est le point de la droite $(d)$ d'abscisse $2$ et $H$ est le projeté orthogonal du point $B$ sur le plan $\mathcal P$ d'équation
\[
x+3z-7=0.
\]

\par\medskip
\noindent\textbf{Affirmation 1.}\quad La droite $(d)$ et l'axe des ordonnées sont deux droites non coplanaires.

\medskip
\noindent\textbf{Affirmation 2.}\quad Le plan passant par $A$ et orthogonal à la droite $(d)$ a pour équation cartésienne :
\[
  x+3z+3=0.
\]

\medskip
\noindent\textbf{Affirmation 3.}\quad Une mesure, exprimée en radian, de l'angle géométrique $\widehat{BAC}$ est $\dfrac{\pi}{6}$.

\medskip
\noindent\textbf{Affirmation 4.}\quad La distance $BH$ est égale à $\dfrac{\sqrt{10}}{2}$.

\newpage
\titreSujet{Sujet 2 -- Plan, projeté orthogonal et tétraèdre}{Polynésie, 18 juin 2025, sujet 2, exercice 3}

L'espace est rapporté à un repère orthonormé $(O;\vec{i},\vec{j},\vec{k})$.
On considère les points suivants :
\[
A(1;3;0),\qquad B(-1;4;5),\qquad C(0;1;0),\qquad D(-2;2;1).
\]

\begin{enumerate}
  \item Montrer que les points $A$, $B$ et $C$ déterminent un plan.
  \item Montrer que le triangle $ABC$ est rectangle en $A$.
  \item Soit $\Delta$ la droite passant par le point $D$ et de vecteur directeur
  \[
  \vec{u}\begin{pmatrix}2\\-1\\1\end{pmatrix}.
  \]
  \begin{enumerate}
    \item Démontrer que la droite $\Delta$ est orthogonale au plan $(ABC)$.
    \item Justifier que le plan $(ABC)$ admet pour équation cartésienne :
    \[
    2x-y+z+1=0.
    \]
    \item Déterminer une représentation paramétrique de la droite $\Delta$.
  \end{enumerate}
  \item On appelle $H$ le point de coordonnées
  \[
  H\left(-\dfrac{2}{3};\dfrac{4}{3};\dfrac{5}{3}\right).
  \]
  Vérifier que $H$ est le projeté orthogonal du point $D$ sur le plan $(ABC)$.
  \item On rappelle que le volume d'un tétraèdre est donné par
  \[
  V=\dfrac13 B\times h,
  \]
  où $B$ est l'aire d'une base du tétraèdre et $h$ est sa hauteur relative à cette base.
  \begin{enumerate}
    \item Montrer que $DH=\dfrac{2\sqrt6}{3}$.
    \item En déduire le volume du tétraèdre $ABCD$.
  \end{enumerate}
  \item On considère la droite $d$ de représentation paramétrique :
  \[
  \begin{cases}
  x=1-2k,\\
  y=-3k,\\
  z=1+k,
  \end{cases}\qquad k\in\mathbb{R}.
  \]
  La droite $d$ et le plan $(ABC)$ sont-ils sécants ou parallèles ?
\end{enumerate}

\newpage
\titreSujet{Sujet 3 -- Drones et plan dans l'espace}{Polynésie, 20 juin 2024, sujet 2, exercice 4}

Une commune décide de remplacer le traditionnel feu d'artifice du 14 juillet par un spectacle de drones lumineux.

Pour le pilotage des drones, l'espace est muni d'un repère orthonormé $(O;\vec{i},\vec{j},\vec{k})$ dont l'unité est la centaine de mètres.
La position de chaque drone est modélisée par un point et chaque drone est envoyé d'un point de départ $D$ de coordonnées $(2;5;1)$.

On souhaite former avec des drones des figures en les positionnant dans un même plan $\mathcal P$.
Trois drones sont positionnés aux points
\[
A(-1;-1;17),\qquad B(4;-2;4),\qquad C(1;-3;7).
\]

\begin{enumerate}
  \item Justifier que les points $A$, $B$ et $C$ ne sont pas alignés.
\end{enumerate}
Dans la suite, on note $\mathcal P$ le plan $(ABC)$ et on considère le vecteur
\[
\vec n\begin{pmatrix}2\\-3\\1\end{pmatrix}.
\]
\begin{enumerate}[resume]
  \item
  \begin{enumerate}
    \item Justifier que $\vec n$ est normal au plan $(ABC)$.
    \item Démontrer qu'une équation cartésienne du plan $\mathcal P$ est
    \[
    2x-3y+z-18=0.
    \]
  \end{enumerate}
  \item Le pilote des drones décide d'envoyer un quatrième drone en prenant comme trajectoire la droite $d$ dont une représentation paramétrique est donnée par
  \[
  d:\begin{cases}
  x=3t+2,\\
  y=t+5,\\
  z=4t+1,
  \end{cases}\qquad t\in\mathbb{R}.
  \]
  \begin{enumerate}
    \item Déterminer un vecteur directeur de la droite $d$.
    \item Afin que ce nouveau drone soit également placé dans le plan $\mathcal P$, déterminer par le calcul les coordonnées du point $E$, intersection de la droite $d$ avec le plan $\mathcal P$.
  \end{enumerate}
  \item Le pilote des drones décide d'envoyer un cinquième drone le long de la droite $\Delta$ qui passe par le point $D$ et qui est perpendiculaire au plan $\mathcal P$. Ce cinquième drone est placé lui aussi dans le plan $\mathcal P$, soit à l'intersection entre la droite $\Delta$ et le plan $\mathcal P$.

  On admet que le point $F(6;-1;3)$ correspond à cet emplacement.

  Démontrer que la distance entre le point de départ $D$ et le plan $\mathcal P$ vaut $2\sqrt{14}$ centaines de mètres.
  \item L'organisatrice du spectacle demande au pilote d'envoyer un nouveau drone dans le plan, peu importe sa position dans le plan, toujours à partir du point $D$.

  Sachant qu'il reste $40$ secondes avant le début du spectacle et que le drone vole en trajectoire rectiligne à $18{,}6$ m.s$^{-1}$, le nouveau drone peut-il arriver à temps ?
\end{enumerate}
\end{document}
